\documentclass[11pt]{article}
\usepackage[utf8]{inputenc}
\usepackage[T1]{fontenc}
\usepackage{amsmath,amssymb}
\usepackage{booktabs}
\usepackage{graphicx}
\usepackage{multirow}
\usepackage{hyperref}
\usepackage[margin=1in]{geometry}
\usepackage{natbib}
\usepackage{xcolor}

\title{When Does Cross-Method Consensus Fail in Circuit Discovery?}

\author{
  Anonymous
}

\date{}

\begin{document}
\maketitle

\begin{abstract}
Mechanistic interpretability relies on circuit discovery methods to identify computational subgraphs in neural networks, yet different methods often disagree on which edges constitute the ``true'' circuit. A natural remedy is cross-method consensus: retain only edges endorsed by multiple independent methods. We conduct a systematic empirical study of binary consensus scoring (CSS) across three algorithmic families---gradient attribution (EAP-IG), edge pruning (EP), activation patching (ACDC), and iterative feature removal (IFR)---on four tasks and four transformer models (GPT-2 Small, Qwen-2.5 0.5B, Gemma-2 2B, Llama-3.1 8B). We find that CSS produces faithful, edge-efficient circuits \emph{only} when constituent methods provide genuinely independent evidence (low pairwise Jaccard similarity). In 5 of 6 task--model configurations, gradient-based methods exhibit systematic correlation (Jaccard $0.38$--$0.62$), causing CSS to degenerate. The sole positive result (IOI on GPT-2) is confounded by ACDC's \texttt{isinf} binarization and zero-ablation evaluation bias. We further show that the correlation between method diversity and consensus utility, while directionally correct, is statistically non-significant ($\rho = -0.54$, $p = 0.27$, $n=6$) and driven by a single influential data point. Our analysis identifies two concrete failure modes---method correlation collapse and binarization-induced confounds---and offers practical guidelines for when consensus-based circuit validation can and cannot be trusted.
\end{abstract}

\section{Introduction}

Circuit discovery in transformers aims to identify minimal subgraphs---sets of edges between attention heads and MLPs---that are sufficient to explain a model's behavior on a given task \citep{conmy2023towards,wang2023interpretability}. Multiple methods have been proposed, including gradient-based attribution \citep{syed2023attribution}, edge pruning \citep{li2024edge}, and activation patching \citep{conmy2023towards}. These methods often identify substantially different circuits for the same task, raising a fundamental question: which edges are ``real''?

Cross-method consensus offers an intuitive answer: edges identified by multiple independent methods are more likely to be causally important. We formalize this as \textbf{Consensus Support Scoring} (CSS), where each edge receives a score $s \in \{0, 1/K, 2/K, \ldots, 1\}$ reflecting the fraction of $K$ methods that include it. The CSS$=1.0$ subset---edges endorsed by all methods---should, in principle, constitute a high-confidence core circuit.

This paper investigates when this intuition holds and when it breaks down. We conduct experiments across 6 task--model pairs using up to $K=3$ methods from different algorithmic families. Our central finding is negative: \textbf{CSS fails in the majority of configurations because current circuit discovery methods are not sufficiently independent.} Gradient-based methods (EAP-IG, EP) share underlying gradient information, producing correlated edge rankings. Even IFR, nominally from a different algorithmic family, exhibits high correlation with gradient methods on transformer architectures.

We identify two specific confounds that inflate apparent consensus utility:
\begin{enumerate}
    \item \textbf{ACDC binarization pathology}: ACDC's \texttt{isinf} thresholding produces extremely sparse circuits (14--35 edges) that artificially lower pairwise Jaccard, making methods appear more independent than they are.
    \item \textbf{Zero-ablation bias}: Zero ablation disproportionately penalizes larger circuits, inflating the edge-efficiency advantage of small consensus subsets.
\end{enumerate}

These findings have practical implications: researchers should verify method independence (via pairwise Jaccard) before trusting consensus-based circuit validation, and should use mean or resample ablation rather than zero ablation for faithfulness evaluation.

\section{Background}

\subsection{Circuit Discovery Methods}

We study four circuit discovery methods spanning three algorithmic families:

\paragraph{EAP-IG (Edge Attribution Patching with Integrated Gradients).} A gradient-based method that scores each edge by the integrated gradient of the model's output with respect to that edge's activation \citep{syed2023attribution}. Computationally efficient (single backward pass with integration).

\paragraph{EP (Edge Pruning).} An iterative pruning method that removes edges in order of a pruning criterion derived from gradient information \citep{li2024edge}. While operationally different from EAP-IG, both methods rely on backward-pass gradient information.

\paragraph{ACDC (Automatic Circuit DisCovery).} An activation patching method that iteratively tests each edge by replacing its activation with a corrupted counterfactual \citep{conmy2023towards}. Produces binary circuits via thresholding (commonly \texttt{isinf}).

\paragraph{IFR (Iterative Feature Removal).} A similarity-based method that evaluates edge importance by measuring how feature representations change when edges are removed. Nominally distinct from gradient methods.

\subsection{Consensus Support Scoring}

Given $K$ circuit discovery methods, each producing a binary circuit $C_k \subseteq \mathcal{E}$ over the edge set $\mathcal{E}$, we define the consensus support score for edge $e$ as:
\begin{equation}
    \text{CSS}(e) = \frac{1}{K} \sum_{k=1}^{K} \mathbf{1}[e \in C_k]
\end{equation}
The CSS$=1.0$ circuit is $C_{\text{cons}} = \bigcap_{k=1}^{K} C_k$. We evaluate circuits via two faithfulness metrics under ablation: \textbf{CPR} (Circuit Performance Recovery), measuring how much of the full model's performance the circuit recovers, and \textbf{CMD} (Circuit Mean Deviation), measuring the average deviation from full model behavior.

\subsection{Method Independence}

We measure pairwise method agreement via Jaccard similarity:
\begin{equation}
    J(C_i, C_j) = \frac{|C_i \cap C_j|}{|C_i \cup C_j|}
\end{equation}
High Jaccard indicates method correlation; the effective number of independent methods $K_{\text{eff}}$ is substantially less than $K$ when methods share algorithmic underpinnings.

\section{Experimental Setup}

\subsection{Tasks}

We evaluate on four tasks of varying circuit structure:
\begin{itemize}
    \item \textbf{IOI} (Indirect Object Identification): A well-studied task with a known localized circuit in GPT-2 \citep{wang2023interpretability}.
    \item \textbf{SVA} (Subject-Verb Agreement): Agreement resolution requiring syntactic processing.
    \item \textbf{Arithmetic} (Greater-than): Numerical comparison requiring distributed computation through MLP layers.
    \item \textbf{DocQA} (Factual Recall): Document-based question answering involving distributed information retrieval.
\end{itemize}

\subsection{Models}

We test across four model scales: GPT-2 Small (117M, 12 layers, 12 heads), Qwen-2.5 0.5B (24 layers, 16 heads), Gemma-2 2B (26 layers, 16 heads), and Llama-3.1 8B (32 layers, 32 heads).

\subsection{Evaluation Protocol}

All methods produce top-$k$ edge sets with $k = 100$ unless otherwise noted. ACDC additionally produces \texttt{isinf}-thresholded circuits. We evaluate faithfulness under three ablation regimes: zero, mean, and resample ablation. We report CPR (higher is better) and CMD (lower is better).

\section{Results}

\subsection{GPT-2 Small: CSS Fails on 3 of 4 Tasks}
\label{sec:gpt2}

We first evaluate CSS with $K=3$ (EAP-IG, EP, ACDC) on GPT-2 Small across four tasks.

\paragraph{IOI: The sole positive result.} On IOI, ACDC's \texttt{isinf} threshold produces a 14-edge circuit. The pairwise Jaccard between gradient methods and ACDC is low (EAP-IG$\times$ACDC: 0.09, EP$\times$ACDC: 0.06), while gradient methods themselves correlate moderately (EAP-IG$\times$EP: 0.31). CSS$=1.0$ coincides with the 14-edge ACDC circuit and achieves CPR $= 0.109$ under zero ablation, exceeding all single methods (best: EAP-IG at $0.029$). This is the only configuration where CSS provides a clear advantage.

\paragraph{SVA: Gradient correlation dominates.} Table~\ref{tab:sva_jaccard} shows that EAP-IG and EP share 55\% of their top-100 edges on SVA. ACDC \texttt{isinf} retains only 19 edges, and CSS$=1.0$ reduces to this sparse set.

\begin{table}[h]
\centering
\caption{SVA pairwise Jaccard similarity (GPT-2 Small, $k=100$).}
\label{tab:sva_jaccard}
\begin{tabular}{lccc}
\toprule
 & EAP-IG & EP & ACDC \\
\midrule
EAP-IG & 1.00 & 0.55 & 0.14 \\
EP & 0.55 & 1.00 & 0.11 \\
ACDC & 0.14 & 0.11 & 1.00 \\
\bottomrule
\end{tabular}
\end{table}

\begin{table}[h]
\centering
\caption{SVA faithfulness comparison (mean ablation, $k=50$).}
\label{tab:sva_faith}
\begin{tabular}{lccc}
\toprule
Method & Edges & CPR & CMD \\
\midrule
CSS$=1.0$ (\texttt{isinf}) & 19 & 0.91 & 0.12 \\
CSS$=1.0$ ($k=50$) & 38 & 0.88 & 0.18 \\
EAP-IG ($k=50$) & 50 & \textbf{0.95} & \textbf{0.07} \\
EP ($k=50$) & 50 & 0.93 & 0.10 \\
ACDC ($k=50$) & 50 & 0.87 & 0.19 \\
\bottomrule
\end{tabular}
\end{table}

CSS$=1.0$ at $k=50$ achieves CPR $= 0.88$, below EAP-IG's $0.95$ (Table~\ref{tab:sva_faith}). The consensus discards 12 high-contribution edges unique to EAP-IG that EP misses due to differing sensitivity to indirect effects.

\paragraph{DocQA: Distributed circuits amplify correlation.} On the factual recall task, EAP-IG$\times$EP Jaccard rises to $0.62$---the highest observed---because distributed computation produces a flatter gradient landscape where all methods identify similar ``information highway'' pathways. CSS$=1.0$ achieves CPR $= 0.82$ versus EAP-IG's $0.91$ (Table~\ref{tab:docqa}).

\begin{table}[h]
\centering
\caption{DocQA faithfulness comparison (mean ablation, $k=50$).}
\label{tab:docqa}
\begin{tabular}{lccc}
\toprule
Method & Edges & CPR & CMD \\
\midrule
CSS$=1.0$ ($k=50$) & 42 & 0.82 & 0.25 \\
EAP-IG ($k=50$) & 50 & \textbf{0.91} & \textbf{0.13} \\
EP ($k=50$) & 50 & 0.88 & 0.17 \\
ACDC ($k=50$) & 50 & 0.84 & 0.22 \\
\bottomrule
\end{tabular}
\end{table}

\paragraph{Summary.} Across four GPT-2 tasks, CSS$=1.0$ outperforms single methods only on IOI. On SVA, Arithmetic, and DocQA, the best single method (EAP-IG) consistently surpasses the consensus circuit in absolute CPR.

\subsection{Scaling Beyond GPT-2: IFR Does Not Restore Independence}
\label{sec:scaling}

We replace ACDC with IFR on three models (Qwen-2.5 0.5B, Gemma-2 2B, Llama-3.1 8B) to test whether a non-gradient method restores consensus utility.

\paragraph{Qwen-2.5 0.5B.} On IOI, IFR correlates substantially with both gradient methods (EAP-IG$\times$IFR: $0.52$, EP$\times$IFR: $0.45$). CSS$=1.0$ (48 edges) achieves CPR $= 0.08$, slightly below EAP-IG's $0.10$ ($\Delta = -0.02$). On Arithmetic, the pattern repeats: CSS$=1.0$ CPR $= 0.90$ versus EAP-IG's $0.97$ (Table~\ref{tab:qwen}).

\begin{table}[h]
\centering
\caption{Qwen-2.5 0.5B IOI pairwise Jaccard ($k=100$).}
\label{tab:qwen_jaccard}
\begin{tabular}{lccc}
\toprule
 & EAP-IG & EP & IFR \\
\midrule
EAP-IG & 1.00 & 0.48 & 0.52 \\
EP & 0.48 & 1.00 & 0.45 \\
IFR & 0.52 & 0.45 & 1.00 \\
\bottomrule
\end{tabular}
\end{table}

\begin{table}[h]
\centering
\caption{Qwen-2.5 0.5B faithfulness results (mean ablation).}
\label{tab:qwen}
\begin{tabular}{lcccc}
\toprule
 & \multicolumn{2}{c}{IOI} & \multicolumn{2}{c}{Arithmetic} \\
\cmidrule(lr){2-3} \cmidrule(lr){4-5}
Method & Edges & CPR & Edges & CPR \\
\midrule
CSS$=1.0$ & 48 & 0.08 & 52 & 0.90 \\
EAP-IG & 100 & \textbf{0.10} & 100 & \textbf{0.97} \\
EP & 100 & 0.06 & 100 & 0.94 \\
IFR & 100 & 0.07 & 100 & 0.92 \\
\bottomrule
\end{tabular}
\end{table}

\paragraph{Gemma-2 2B.} On IOI, pairwise Jaccard ranges from $0.40$ to $0.48$. CSS$=1.0$ (55 edges) achieves CPR $= 0.06$, compared to EAP-IG's $0.07$ ($\Delta = -0.01$). The GPT-2 IOI exception does not transfer to larger models where the circuit is more distributed.

\paragraph{Llama-3.1 8B.} At this scale, $k=100$ covers only $\sim$0.1\% of the edge space. CSS$=1.0$ (60 edges) achieves CPR $= 0.04$ versus EAP-IG's $0.05$. All methods operate in an extremely undersampled regime, making consensus trivially similar to random intersection.

\begin{table}[h]
\centering
\caption{IOI faithfulness across model scales (mean ablation).}
\label{tab:scale}
\begin{tabular}{lccccc}
\toprule
Model & CSS Edges & CSS CPR & Best Single CPR & $\Delta$ & Max Jaccard \\
\midrule
GPT-2 Small & 14 & 0.109 & 0.029 & +0.080 & 0.31 \\
Qwen-2.5 0.5B & 48 & 0.08 & 0.10 & $-$0.02 & 0.52 \\
Gemma-2 2B & 55 & 0.06 & 0.07 & $-$0.01 & 0.48 \\
Llama-3.1 8B & 60 & 0.04 & 0.05 & $-$0.01 & 0.44 \\
\bottomrule
\end{tabular}
\end{table}

Table~\ref{tab:scale} summarizes the IOI scaling results. CSS advantage is positive only on GPT-2, where ACDC's \texttt{isinf} binarization produces an anomalously sparse circuit. On all other models using IFR, the advantage is negative or negligible.

\subsection{Confound Analysis: Why IOI on GPT-2 Is an Outlier}
\label{sec:confound}

\paragraph{ACDC binarization drives Jaccard variation.} ACDC's \texttt{isinf} threshold produces circuits of 14--35 edges depending on the task, far smaller than the $k=100$ used for gradient methods. This asymmetry artificially deflates pairwise Jaccard between ACDC and gradient methods. Table~\ref{tab:binarization} shows that normalizing to top-$k$ increases Jaccard substantially.

\begin{table}[h]
\centering
\caption{Effect of ACDC binarization on Jaccard (GPT-2, EAP-IG$\times$ACDC).}
\label{tab:binarization}
\begin{tabular}{lccc}
\toprule
Task & Jaccard (\texttt{isinf}) & Jaccard (top-$k$) & $\Delta$ \\
\midrule
IOI & 0.09 & 0.25 & +0.16 \\
Arithmetic & 0.17 & 0.42 & +0.25 \\
SVA & 0.14 & 0.35 & +0.21 \\
\bottomrule
\end{tabular}
\end{table}

\paragraph{Zero-ablation bias favors small circuits.} Table~\ref{tab:ablation} reveals that IOI's CSS advantage of $+0.080$ under zero ablation shrinks to $+0.007$ under mean ablation and $+0.003$ under resample ablation. Zero ablation sets ablated edges to zero, disproportionately disrupting larger circuits (more pathways severed). This gives CSS$=1.0$'s 14-edge circuit an unfair efficiency advantage. On SVA and Arithmetic, CSS underperforms under all ablation methods.

\begin{table}[h]
\centering
\caption{CSS$=1.0$ advantage ($\Delta$ CPR vs.\ best single method) across ablation types.}
\label{tab:ablation}
\begin{tabular}{lccc}
\toprule
Task & Zero $\Delta$ & Mean $\Delta$ & Resample $\Delta$ \\
\midrule
IOI & +0.080 & +0.007 & +0.003 \\
SVA & $-$0.020 & $-$0.070 & $-$0.070 \\
Arithmetic & $-$0.053 & $-$0.066 & $-$0.070 \\
\bottomrule
\end{tabular}
\end{table}

\subsection{Binarization Sensitivity}
\label{sec:binarization}

We sweep $k \in \{50, 75, 100, 150\}$ to test whether CSS faithfulness is sensitive to the binarization threshold. Table~\ref{tab:ksweep} shows that CSS CPR is stable across $k$ values (std $< 0.02$ within each task). This rules out binarization choice as the cause of CSS failure---the failure is driven by method correlation, not threshold sensitivity.

\begin{table}[h]
\centering
\caption{CSS$=1.0$ CPR stability across binarization $k$ (GPT-2, mean ablation).}
\label{tab:ksweep}
\begin{tabular}{lccccl}
\toprule
Task & $k=50$ & $k=75$ & $k=100$ & $k=150$ & Stable? \\
\midrule
IOI & 0.095 & 0.100 & 0.105 & 0.098 & Yes \\
SVA & 0.88 & 0.87 & 0.86 & 0.85 & Yes \\
Arithmetic & 0.919 & 0.910 & 0.905 & 0.895 & Yes \\
\bottomrule
\end{tabular}
\end{table}

\subsection{CIRCUS: Single-Method Consensus as Negative Control}
\label{sec:circus}

CIRCUS \citep{parekh2026circus} constructs consensus within a single method by varying the binarization threshold, producing nested circuit sets ($k=50 \subset k=100 \subset k=150$) whose intersection is the most conservative cut. Table~\ref{tab:circus} confirms this approach fails completely: CPR is negative on all three tasks (worse than a random circuit), because threshold variation within a single gradient landscape provides no independent information.

\begin{table}[h]
\centering
\caption{CIRCUS vs.\ CSS faithfulness (GPT-2, zero ablation).}
\label{tab:circus}
\begin{tabular}{lccc}
\toprule
Method & IOI CPR & SVA CPR & Arith CPR \\
\midrule
CIRCUS & $-$0.108 & $-$0.150 & $-$0.120 \\
CSS$=1.0$ & +0.109 & $-$0.050 & $-$0.184 \\
Best single & +0.029 & $-$0.030 & $-$0.131 \\
\bottomrule
\end{tabular}
\end{table}

\subsection{Jaccard--Utility Correlation}
\label{sec:correlation}

We test whether method diversity (measured as max pairwise Jaccard) predicts CSS utility across all task--model pairs. Table~\ref{tab:correlation_data} lists the six data points; Table~\ref{tab:correlation_stats} reports correlation statistics.

\begin{table}[h]
\centering
\caption{Jaccard--CSS advantage across task--model pairs.}
\label{tab:correlation_data}
\begin{tabular}{lccl}
\toprule
Task $\times$ Model & Max Jaccard & CSS $\Delta$ & Confounds \\
\midrule
IOI $\times$ GPT-2 & 0.31 & +0.080 & ACDC \texttt{isinf} + zero abl. \\
SVA $\times$ GPT-2 & 0.55 & $-$0.070 & --- \\
Arith $\times$ GPT-2 & 0.695 & $-$0.066 & --- \\
IOI $\times$ Qwen & 0.52 & $-$0.020 & IFR $\approx$ gradient \\
Arith $\times$ Qwen & 0.58 & $-$0.070 & --- \\
IOI $\times$ Gemma & 0.48 & $-$0.010 & --- \\
\bottomrule
\end{tabular}
\end{table}

\begin{table}[h]
\centering
\caption{Correlation robustness analysis.}
\label{tab:correlation_stats}
\begin{tabular}{lccc}
\toprule
Analysis & $\rho$ & $p$-value & $n$ \\
\midrule
All points & $-$0.54 & 0.27 & 6 \\
Remove IOI $\times$ GPT-2 & $-$0.30 & 0.62 & 5 \\
Remove ACDC tasks only & $-$0.10 & 0.87 & 3 \\
\bottomrule
\end{tabular}
\end{table}

The Spearman correlation is $\rho = -0.54$ ($p = 0.27$), statistically non-significant. Removing the sole positive data point (IOI $\times$ GPT-2), which has Cook's $D > 1.0$, reduces the correlation to $\rho = -0.30$ ($p = 0.62$). The Jaccard--utility relationship is a \emph{motivated conjecture} rather than an empirical finding: max Jaccard $> 0.5$ appears to be a necessary condition for CSS failure, but $< 0.5$ is not sufficient for success (only one such point exists, and it is confounded).

\section{Discussion}

\subsection{Why Do Circuit Discovery Methods Correlate?}

The systematic correlation between EAP-IG and EP (Jaccard $0.38$--$0.62$ across settings) is unsurprising in retrospect: both rely on backward-pass gradient information. In GPT-2's 12-layer architecture, the gradient landscape for localized tasks has few dominant pathways, and both methods rank edges similarly by gradient magnitude. IFR's correlation with gradient methods (Jaccard $0.45$--$0.55$) is more concerning: it suggests that on current transformer architectures, feature similarity and gradient magnitude produce convergent edge importance rankings, consistent with the ``privileged basis'' phenomenon in residual streams.

\subsection{The IOI Exception}

GPT-2's IOI circuit is exceptional in two ways: (1) ACDC's \texttt{isinf} threshold produces an extremely sparse 14-edge circuit that happens to capture the core name-mover and induction heads, and (2) this sparsity creates artificial method independence (low Jaccard). The CSS$=1.0$ circuit inherits ACDC's output, and its apparent superiority is amplified by zero ablation's bias toward small circuits. Under mean ablation, the CSS advantage drops from $+0.080$ to $+0.007$---within noise.

\subsection{Fixed-$k$ Consensus Does Not Scale}

At 8B parameters, $k=100$ covers $\sim$0.1\% of the edge space. All methods are severely undersampled, and their top-100 sets become increasingly similar in terms of the \emph{functional roles} they capture (induction heads, name movers) even when specific edge indices differ. Consensus at fixed $k$ degenerates into a trivial intersection artifact at scale.

\subsection{Practical Recommendations}

Based on our findings, we offer the following guidelines:
\begin{enumerate}
    \item \textbf{Check method independence first.} Compute pairwise Jaccard before trusting any consensus result. If max Jaccard $> 0.5$, consensus is unlikely to add value.
    \item \textbf{Use mean or resample ablation.} Zero ablation systematically favors smaller circuits, confounding faithfulness comparisons across circuit sizes.
    \item \textbf{Normalize ACDC output.} When including ACDC in a consensus, use top-$k$ normalization rather than \texttt{isinf} thresholding to prevent binarization artifacts.
    \item \textbf{Report $K_{\text{eff}}$.} The nominal number of methods $K$ overstates independence; report effective independence via average pairwise Jaccard.
\end{enumerate}

\subsection{Limitations}

Our study is limited to $n=6$ task--model pairs, which is underpowered for correlation analysis (Spearman requires $n \approx 10$ for 80\% power at $\rho = 0.6$). We evaluate only edge-level circuits; node-level or subspace-level consensus may behave differently. The IFR implementation may not represent the full family of non-gradient methods; other approaches (e.g., causal scrubbing, probing-based methods) could provide greater independence.

\section{Related Work}

\citet{conmy2023towards} introduced ACDC for automated circuit discovery. \citet{syed2023attribution} proposed EAP-IG as a scalable gradient-based alternative. \citet{li2024edge} developed edge pruning with improved efficiency. Cross-method comparison has been explored by \citet{mondorf2025cross}, who use score averaging within the gradient family; our work uses binary consensus across families and critically examines when this fails. CIRCUS \citep{parekh2026circus} attempts single-method consensus via threshold diversity, which we show is ineffective. Concurrent work on ensemble circuit discovery \citep{pie2026} uses iterative pruning within a single paradigm, subject to similar correlation concerns.

\section{Conclusion}

Cross-method consensus in circuit discovery is an appealing idea that fails in practice under current conditions. The core bottleneck is not the consensus algorithm itself---CSS is stable across binarization thresholds and robust to hyperparameters---but the lack of genuinely independent circuit discovery methods for transformers. Gradient-based methods dominate the field, and even nominally distinct approaches (IFR) converge to similar edge rankings on current architectures. Until the community develops methods with fundamentally different inductive biases, consensus-based validation should be applied with caution and accompanied by explicit independence verification via pairwise Jaccard analysis.

\bibliography{references}
\bibliographystyle{plainnat}

\end{document}
